% !TEX root = ../jalase10.tex
%% ============================================================
%%  جلسهٔ ۱۰ — تابعی از یک متغیر تصادفی: روش CDF، تبدیل خطی،
%%  فرمول تغییر متغیر، رادیکال و نمایی، روش CDF وارون
%% ============================================================
\chapter{تابعی از یک متغیر تصادفی}
\label{ch:s10}

\sessioncard{دوشنبه ۴ آبان}
{توزیعِ تابعی از یک متغیر تصادفی؛ روش \lr{CDF} و فرمول تغییر متغیر؛ تبدیل خطی متغیر گاوسی؛ ساختن توزیع‌های دلخواه از توزیع یکنواخت}
{\lr{Pishro-Nik} بخش‌های \lr{3.2.3} و \lr{4.1.3} \,؛\, \lr{Ross} بخش‌های \lr{5.7} و \lr{10.2.1} \,؛\, \lr{Bertsekas} بخش \lr{4.1}}
{\begin{itemize}
  \item اگر توزیع $X$ را بدانید، \lr{CDF} و \lr{PDF} متغیر $Y=g(X)$ را پیدا کنید؛
  \item در تبدیل خطی $aX+b$ حواستان به علامت $a$ باشد و ثابت کنید تبدیل خطی گاوسی، گاوسی است؛
  \item توزیع $\sqrt X$ و $e^X$ را برای توزیع‌های آشنا حساب کنید؛
  \item فقط با یک مولد عدد یکنواخت و \lr{CDF} وارون، از هر توزیعی که بخواهید نمونه بسازید.
\end{itemize}}

دمای یک اتاق را با سنسوری اندازه می‌گیریم که کمی خطا دارد؛ پس عددی که می‌خوانیم، $X$، تصادفی است و فرض کنید توزیعش را هم می‌دانیم. نرم‌افزاری که این عدد را می‌گیرد آن را به فارنهایت می‌برد: $Y=1.8X+32$. توزیع $Y$ چیست؟ یا $X$ توانِ سیگنالی است که به آنتن می‌رسد و ما دامنه‌اش، $\sqrt X$، را لازم داریم. یا $X$ لگاریتمِ بازده یک سهم است و قیمت، $e^X$ است.

سؤال در همهٔ این‌ها یکی است: توزیع $X$ را داریم، توزیع $g(X)$ چیست؟ جوابش هم یک جمله است: هر سؤالی دربارهٔ $Y$ را به سؤالی دربارهٔ $X$ ترجمه کنید. بقیهٔ جلسه جزئیاتِ همین ترجمه است؛ جزئیاتی که اگر حواسمان نباشد، مثلاً به علامت یک ضریب، جواب را خراب می‌کنند.

آخر جلسه همین ایده را برعکس به کار می‌بریم. تابع \lr{\texttt{random()}} در تقریباً هر زبان برنامه‌نویسی فقط عددی یکنواخت بین ۰ و ۱ می‌دهد. ولی اگر بخواهید صف یک بانک را شبیه‌سازی کنید، فاصلهٔ ورود مشتری‌ها باید نمایی باشد، نه یکنواخت. می‌بینیم که با انتخابِ یک تابعِ درست، از همان عدد یکنواخت هر توزیعی که بخواهیم بیرون می‌آید.

%% ------------------------------------------------------------
\section{تابعِ یک متغیر تصادفی هم متغیر تصادفی است}
\label{sec:functions}
%% ------------------------------------------------------------

متغیر تصادفی $X$ تابعی از $\Omega$ به $\R$ است (جلسهٔ ۵). اگر $g$ یک تابع معمولی از $\R$ به $\R$ باشد، می‌شود این دو را پشت سر هم اجرا کرد:
\[ \omega\ \longmapsto\ X(\omega)\ \longmapsto\ g\bigl(X(\omega)\bigr). \]
حاصل باز هم تابعی از $\Omega$ به $\R$ است؛ پس $Y=g(X)$ خودش متغیر تصادفی است و توزیع خودش را دارد. $X^2$، $|X|$، $3X+2$، $e^X$ و $\ind\{X>5\}$ همه متغیر تصادفی‌اند.

راحت‌ترین تصویر این است: آزمایش انجام می‌شود، عددِ $X$ بیرون می‌آید و ما آن را از دستگاهی به نام $g$ رد می‌کنیم. تصادف فقط در $X$ است و $g$ چیزی به آن اضافه نمی‌کند. برای همین هر سؤالی دربارهٔ $Y$ در اصل سؤالی دربارهٔ $X$ است:
\[ \{Y\in B\}=\bigl\{\omega:\ g(X(\omega))\in B\bigr\}=\bigl\{X\in g^{-1}(B)\bigr\}, \]
که در آن $g^{-1}(B)=\{x:\ g(x)\in B\}$ یعنی همهٔ $x$هایی که $g$ آن‌ها را به داخل $B$ می‌فرستد. این $g^{-1}$ لزوماً «تابع وارون» نیست؛ یک مجموعه است و برای هر $g$ای معنی دارد، حتی وقتی $g$ یک‌به‌یک نباشد.

\begin{keyformula}
\[ \P(Y\in B)=\P\bigl(X\in g^{-1}(B)\bigr) \]
\end{keyformula}

هر چه در این جلسه می‌آید، استفادهٔ پشت سر هم از همین تساوی است.

\subsection{حالت گسسته}

اگر $X$ گسسته باشد، $Y$ هم گسسته است و کار سرراست است. برای هر مقدار $y$ ببینید کدام $x$ها به $y$ می‌روند و احتمال‌هایشان را جمع کنید:
\[ p_Y(y)=\sum_{x:\ g(x)=y}p_X(x). \]
اگر $g$ یک‌به‌یک باشد، هر $y$ فقط از یک $x$ می‌آید و عملاً فقط اسم مقدارها عوض می‌شود. جمع وقتی پیش می‌آید که چند $x$ روی یک $y$ بیفتند.

\begin{example}[label=ex:func-disc]{تابعی از یک متغیر گسسته}
$X$ مقادیر $-1$، $0$ و $1$ را هر کدام با احتمال $\tfrac13$ می‌گیرد. توزیع $Y=X^2$ را بیابید.
\begin{solution}
$Y$ فقط $0$ یا $1$ می‌شود. $\{Y=0\}=\{X=0\}$، ولی $\{Y=1\}=\{X=-1\}\cup\{X=1\}$؛ پس
\[ p_Y(0)=\tfrac13,\qquad p_Y(1)=\tfrac13+\tfrac13=\tfrac23. \]
سه مقدارِ هم‌احتمال به دو مقدارِ نابرابر تبدیل شدند، چون $-1$ و $1$ زیر تابع $x^2$ روی هم افتادند.
\end{solution}
\end{example}

%% ------------------------------------------------------------
\section{حالت پیوسته: از \lr{CDF} شروع کنید}
\label{sec:cdf-method}
%% ------------------------------------------------------------

برای $X$ پیوسته فرمول بالا به کار نمی‌آید. احتمال هر نقطه صفر است و از جمع کردن صفرها چیزی درنمی‌آید. چگالی را هم نمی‌شود همین‌طور از جایی به جای دیگر برد، چون چگالی احتمال نیست (جلسهٔ ۷). اما $\{Y\le y\}$ یک \emph{پیشامد} است و احتمالش را می‌شود از روی $X$ حساب کرد. پس از \lr{CDF} شروع می‌کنیم و چگالی را آخر کار با مشتق‌گیری به دست می‌آوریم.

\begin{keyformula}
\begin{enumerate}[leftmargin=1.4em, itemsep=4pt]
  \item \textbf{برد $Y$:} ببینید $Y$ چه مقدارهایی می‌تواند بگیرد. بیرون از این برد $f_Y=0$ است.
  \item \textbf{ترجمه:} برای $y$های داخل برد، نامساویِ $g(X)\le y$ را بر حسب $X$ حل کنید و $F_Y(y)=\P(g(X)\le y)$ را بر حسب $F_X$ بنویسید.
  \item \textbf{مشتق:} $f_Y(y)=F_Y'(y)$.
\end{enumerate}
\end{keyformula}

فکرِ اصلی در قدم دوم است. قدم سوم معمولاً فقط قاعدهٔ زنجیره‌ای است. قدم اول را هم جدی بگیرید: بیشترِ اشتباه‌ها در این مبحث از این است که یادمان رفته $\sqrt X$ منفی نمی‌شود یا $e^X$ همیشه مثبت است.

%% ------------------------------------------------------------
\section{تبدیل خطی: $Y=aX+b$}
\label{sec:linear}
%% ------------------------------------------------------------

از ساده‌ترین تابع شروع می‌کنیم. $X$ پیوسته با چگالی $f_X$ است، $a\ne0$ و $b$ دو عدد ثابت‌اند و $Y=aX+b$. در قدم دوم باید نامساوی $aX+b\le y$ را برای $X$ حل کنیم، و همین‌جا علامت $a$ مهم می‌شود: \key{تقسیم بر عدد منفی جهت نامساوی را عوض می‌کند.}

\textbf{اگر $a>0$:}
\[ F_Y(y)=\P(aX+b\le y)=\P\Bigl(X\le\frac{y-b}{a}\Bigr)=F_X\Bigl(\frac{y-b}{a}\Bigr),\qquad f_Y(y)=\frac1a\,f_X\Bigl(\frac{y-b}{a}\Bigr). \]

\textbf{اگر $a<0$:}
\[ F_Y(y)=\P(aX+b\le y)=\P\Bigl(X\ge\frac{y-b}{a}\Bigr)=1-F_X\Bigl(\frac{y-b}{a}\Bigr),\qquad f_Y(y)=-\frac1a\,f_X\Bigl(\frac{y-b}{a}\Bigr). \]
(برای متغیر پیوسته $\P(X\ge t)=\P(X>t)=1-F_X(t)$.) منفیِ پشت $\frac1a$ از مشتقِ «$1-\dots$» آمده و چون $a<0$ است، $-\frac1a$ مثبت است؛ پس چگالی مثبت ماند، همان‌طور که باید. در هر دو حالت ضریب جلوی $f_X$ همان $\frac1{|a|}$ است:

\begin{proposition}[label=prop:linear]{چگالی تبدیل خطی}
اگر $X$ پیوسته باشد و $Y=aX+b$ با $a\ne0$، آنگاه
\[ f_Y(y)=\frac{1}{|a|}\,f_X\Bigl(\frac{y-b}{a}\Bigr). \]
\lr{CDF} اما دو شکل دارد: برای $a>0$، $F_Y(y)=F_X(\frac{y-b}{a})$ و برای $a<0$، $F_Y(y)=1-F_X(\frac{y-b}{a})$.
\end{proposition}

هر دو اثرِ این فرمول را می‌شود با چشم دید. ضریب $\frac1{|a|}$ از کش آمدن می‌آید: اگر $|a|=3$ باشد همهٔ فاصله‌ها سه برابر می‌شوند، همان احتمال روی بازه‌ای سه برابر پهن‌تر پخش می‌شود و ارتفاع چگالی باید یک‌سوم شود تا مساحت ۱ بماند. منفی بودن $a$ هم نمودار را آینه می‌کند، چون حالا مقدارهای بزرگ $X$ به مقدارهای کوچک $Y$ می‌روند. $b$ در هیچ‌کدام نقشی ندارد و فقط نمودار را جابه‌جا می‌کند.

\begin{warning}
فرض کنید برای $a<0$ همان فرمولِ حالت $a>0$ را بنویسیم: $F_Y(y)=F_X(\frac{y-b}{a})$. وقتی $y\to+\infty$، عبارت $\frac{y-b}{a}$ به $-\infty$ می‌رود و این «\lr{CDF}» به صفر میل می‌کند؛ یعنی تابعی ساخته‌ایم که نزولی است، و هیچ \lr{CDF}ای نزولی نیست. مشتقش هم چگالی منفی می‌دهد. پس آخر کار همیشه دو چیز را نگاه کنید: \lr{CDF} صعودی است و از ۰ به ۱ می‌رود، و چگالی هیچ‌جا منفی نیست. همین دو چک بیشترِ اشتباه‌های علامت را می‌گیرند.
\end{warning}

\subsection{تبدیل خطی متغیر گاوسی}

در جلسهٔ ۸ ثابت کردیم تبدیل خطی گاوسی، گاوسی است، ولی فقط برای ضریب مثبت؛ درباره حالت  منفی چیزی نگفتیم . حالا با گزارهٔ \ref{prop:linear} هر دو حالت را یک‌جا می‌بینیم.

\begin{theorem}[label=thm:gauss-lin]{تبدیل خطی گاوسی}
اگر $X\sim\mathcal N(\mu,\sigma^2)$ و $a\ne0$ و $b$ دو عدد ثابت باشند، آنگاه
\[ Y=aX+b\ \sim\ \mathcal N\bigl(a\mu+b,\ a^2\sigma^2\bigr). \]
\end{theorem}

\begin{proof}
طبق گزارهٔ \ref{prop:linear}، برای هر علامتی از $a$:
\[ f_Y(y)=\frac1{|a|}\cdot\frac{1}{\sigma\sqrt{2\pi}}\exp\Bigl(-\frac{(\frac{y-b}{a}-\mu)^2}{2\sigma^2}\Bigr). \]
صورت و مخرجِ کسرِ داخل توان را در $a^2$ ضرب می‌کنیم:
\[ \frac{\bigl(\frac{y-b}{a}-\mu\bigr)^2}{2\sigma^2}=\frac{(y-b-a\mu)^2}{2a^2\sigma^2}=\frac{\bigl(y-(a\mu+b)\bigr)^2}{2(|a|\sigma)^2}. \]
پس
\[ f_Y(y)=\frac{1}{(|a|\sigma)\sqrt{2\pi}}\exp\Bigl(-\frac{\bigl(y-(a\mu+b)\bigr)^2}{2(|a|\sigma)^2}\Bigr), \]
که یک توزیع  گاوسی با مرکز $a\mu+b$ و پهنای $|a|\sigma$ است.
\end{proof}

%% ------------------------------------------------------------
\section{فرمول کلی برای تابع‌های یکنوا}
\label{sec:monotone}
%% ------------------------------------------------------------

استدلالی که برای $aX+b$ کردیم فقط به یک ویژگیِ این تابع تکیه داشت: \emph{یکنوا} بود، یعنی یا همه‌جا صعودی یا همه‌جا نزولی. همین استدلال برای هر تابع یکنوای دیگری هم کار می‌کند.

فرض کنید $g$ روی برد $X$ اکیداً یکنوا و مشتق‌پذیر باشد و $h=g^{-1}$ تابع وارونش باشد؛ یعنی اگر $y=g(x)$ آنگاه  $x=h(y)$. اگر $g$ صعودی باشد، می‌توان نتیجه گرفت  اگر $g(X)\le y$  آنگاه $X\le h(y)$ است؛ اگر نزولی باشد، آنگاه $X\ge h(y)$:
\[ \text{\rl{صعودی:}}\ \ F_Y(y)=F_X\bigl(h(y)\bigr),\qquad\qquad \text{\rl{نزولی:}}\ \ F_Y(y)=1-F_X\bigl(h(y)\bigr). \]
با مشتق‌گیری، در حالت صعودی $f_Y(y)=f_X(h(y))\,h'(y)$ و در حالت نزولی $f_Y(y)=-f_X(h(y))\,h'(y)$. در حالت نزولی خودِ $h'$ منفی است، پس باز هم جواب مثبت است. دو حالت در یک فرمول جمع می‌شوند:

\begin{theorem}[label=thm:change-var]{فرمول تغییر متغیر}
اگر $g$ روی برد $X$ اکیداً یکنوا و مشتق‌پذیر باشد و $h=g^{-1}$، آنگاه برای $y$های داخل برد $Y=g(X)$
\[ f_Y(y)=f_X\bigl(h(y)\bigr)\,\bigl|h'(y)\bigr|, \]
و بیرون از این برد $f_Y(y)=0$.
\end{theorem}

\textbf{چرا $|h'(y)|$؟} یک بازهٔ خیلی کوچک $[x,\,x+dx]$ را در نظر بگیرید. $g$ آن را به بازهٔ کوچک $[y,\,y+dy]$ می‌برد که $y=g(x)$ و طولش $|dy|=|g'(x)|\,|dx|$ است. «$X$ در بازهٔ اول است» و «$Y$ در بازهٔ دوم است» یک پیشامدند، پس احتمالشان برابر است:
\begin{keyformula}
\[ f_Y(y)\,|dy|\ =\ f_X(x)\,|dx| \]
\end{keyformula}
یعنی $f_Y(y)=f_X(x)\,\bigl|\frac{dx}{dy}\bigr|$، و $\frac{dx}{dy}$ همان $h'(y)$ است. احتمال جابه‌جا می‌شود ولی کم و زیاد نمی‌شود. مثل خمیری که روی میز پهن می‌کنید: مقدار خمیر ثابت است، ولی هر جا بیشتر کشیده شود نازک‌تر می‌شود. جایی که $g$ تند است ($|g'|$ بزرگ) چگالی $Y$ رقیق می‌شود و جایی که $g$ کُند است چگالی غلیظ می‌شود.

\begin{warning}
ضریبِ $|h'(y)|$ را جا نیندازید. اگر $X\sim\mathrm{Unif}(0,1)$ و $Y=2X$، کسی که فقط $x$ را با $\frac y2$ عوض کند می‌نویسد $f_Y(y)=f_X(\frac y2)=1$ روی $[0,2]$؛ ولی مساحت زیر این «چگالی» ۲ است! جواب درست $\frac12$ روی $[0,2]$ است ($Y\sim\mathrm{Unif}(0,2)$). اگر جوابتان روی کل برد انتگرال ۱ ندارد، به احتمال زیاد همین ضریب را فراموش کرده‌اید.
\end{warning}

این فرمول فقط برای تابع یکنوا درست است. اگر $g$ یکنوا نباشد، روش \lr{CDF} هنوز کار می‌کند؛ نمونه‌اش را در بخش \ref{sec:nonmono} می‌بینیم.

%% ------------------------------------------------------------
\section{دو تابع پرکاربرد: رادیکال و نمایی}
\label{sec:sqrt-exp}
%% ------------------------------------------------------------

\subsection{رادیکال: $Y=\sqrt X$}

برای اینکه $\sqrt X$ تعریف شود فرض می‌کنیم $X\ge0$. پس $Y\ge0$ هم هست و برای $y<0$ داریم $F_Y(y)=0$. برای $y\ge0$، چون رادیکال صعودی است،
\[ F_Y(y)=\P(\sqrt X\le y)=\P(X\le y^2)=F_X(y^2),\qquad f_Y(y)=2y\,f_X(y^2). \]
فرمول تغییر متغیر هم همین را می‌دهد: $h(y)=y^2$ و $h'(y)=2y$.

\begin{example}[label=ex:sqrt]{رادیکالِ یک متغیر یکنواخت}
اگر $X\sim\mathrm{Unif}(0,1)$، توزیع $Y=\sqrt X$ را بیابید.
\begin{solution}
برد $Y$ بازهٔ $[0,1]$ است. برای $0\le y\le1$:
\[ F_Y(y)=F_X(y^2)=y^2,\qquad f_Y(y)=2y. \]
$X$ یکنواخت بود ولی $\sqrt X$ نیست و بیشتر نزدیکِ ۱ است. دلیلش این است که رادیکال عددهای کوچک را بزرگ می‌کند: $\sqrt{0.01}=0.1$ و $\sqrt{0.25}=0.5$. احتمالی که نزدیک صفر بود به سمت راست کشیده می‌شود. مثلاً $\P(Y\le\tfrac12)=\P(X\le\tfrac14)=\tfrac14$، و میانهٔ $Y$ (نقطه‌ای که $F_Y$ در آن $\frac12$ می‌شود) $\sqrt{0.5}\approx0.71$ است. این $2y$ را آخر جلسه دوباره می‌بینیم.
\end{solution}
\end{example}

\begin{example}[label=ex:rayleigh]{از توان به دامنه: توزیع ریلی}
در مخابرات بی‌سیم، وقتی سیگنال از مسیرهای زیادی (بازتاب از دیوارها و ساختمان‌ها) به گیرنده می‌رسد، توانِ دریافتی را معمولاً نمایی مدل می‌کنند: $X\sim\mathrm{Exp}(\lambda)$. توزیع دامنه، $Y=\sqrt X$، را بیابید.
\begin{solution}
برای $y\ge0$:
\[ F_Y(y)=F_X(y^2)=1-e^{-\lambda y^2},\qquad f_Y(y)=2y\cdot\lambda e^{-\lambda y^2}=2\lambda y\,e^{-\lambda y^2}. \]
این توزیع \term{ریلی}{Rayleigh} نام دارد و در مخابرات به این پدیده \emph{محوشدگیِ ریلی} می‌گویند. (با $\lambda=\frac{1}{2s^2}$ به شکلی می‌رسیم که معمولاً در کتاب‌های مخابرات می‌بینید: $f_Y(y)=\frac{y}{s^2}e^{-y^2/(2s^2)}$.)

\end{solution}
\end{example}

\subsection{نمایی: $Y=e^X$}

$e^X$ همیشه مثبت است؛ پس برای $y\le0$، $F_Y(y)=0$. تابع $e^x$ صعودی است و وارونش $\ln$؛ برای $y>0$:
\[ F_Y(y)=\P(e^X\le y)=\P(X\le\ln y)=F_X(\ln y),\qquad f_Y(y)=\frac1y\,f_X(\ln y). \]

\begin{example}[label=ex:lognormal]{نمایِ یک گاوسی: توزیع لگ‌نرمال}
اگر $X\sim\mathcal N(\mu,\sigma^2)$، توزیع $Y=e^X$ را بیابید. برای $\mu=0$ و $\sigma=1$، احتمال‌های $\P(Y\le1)$، $\P(Y>e)$ و $\P(Y>e^2)$ را حساب کنید.
\begin{solution}
با فرمول بالا، برای $y>0$:
\[ F_Y(y)=\Phi\Bigl(\frac{\ln y-\mu}{\sigma}\Bigr),\qquad f_Y(y)=\frac{1}{y\,\sigma\sqrt{2\pi}}\exp\Bigl(-\frac{(\ln y-\mu)^2}{2\sigma^2}\Bigr). \]
به این توزیع \term{لگ‌نرمال}{log-normal} می‌گویند، چون لگاریتمش نرمال (گاوسی) است. برای $\mu=0$ و $\sigma=1$:
\[ \P(Y\le1)=\P(X\le0)=\tfrac12,\qquad \P(Y>e)=\P(X>1)\approx0.159, \]
\[ \P(Y>e^2)=\P(X>2)\approx0.023. \]
نصفِ احتمال در بازهٔ کوتاهِ $(0,1)$ است و نصف دیگر از ۱ تا بی‌نهایت پخش شده. $X$ متقارن بود ولی $Y$ کاملاً کج است: $e^x$ همهٔ اعداد منفی را در $(0,1)$ فشرده می‌کند و اعداد مثبت را تا بی‌نهایت می‌کشد.
\end{solution}
\end{example}

\begin{figblock}
\begin{tikzpicture}[x=1.5cm, y=3.2cm]
  \draw[-{Stealth}, graytext] (-0.1,0) -- (5.4,0) node[right] {$y$};
  \draw[-{Stealth}, graytext] (0,-0.03) -- (0,1.1) node[above] {$f_Y(y)$};
  \draw[graytext, dashed] (1,0) -- (1,1.0);
  \node[above, font=\scriptsize, graytext] at (1,1.0) {\rl{میانه}};
  \draw[setA, very thick, domain=0.02:5, samples=150] plot (\x,{exp(-ln(\x)*ln(\x)/2)/(\x*sqrt(2*pi))});
  \draw[accent, very thick, domain=0.06:5, samples=150] plot (\x,{exp(-ln(\x)*ln(\x)/0.5)/(\x*0.5*sqrt(2*pi))});
  \foreach \tk in {1,2,3,4,5} { \draw[graytext] (\tk,0) -- (\tk,-0.02); \node[below, font=\scriptsize] at (\tk,-0.02) {$\tk$}; }
  \foreach \tk in {0.5,1} { \draw[graytext] (0,\tk) -- (-0.05,\tk); \node[left, font=\scriptsize] at (-0.05,\tk) {$\tk$}; }
  \node[setA!80!black, font=\small, anchor=west] at (2.3,0.25) {$\sigma=1$};
  \node[accent!80!black, font=\small, anchor=west] at (1.5,0.75) {$\sigma=\tfrac12$};
\end{tikzpicture}

\figcaption[fig:lognormal]{چگالی لگ‌نرمال با $\mu=0$. برای هر $\sigma$ میانه در $y=1$ است (چون $\P(X\le0)=\frac12$)، ولی دُمِ راست بسیار بلندتر از دُمِ چپ است. هرچه $\sigma$ بزرگ‌تر باشد این کجی بیشتر است.}
\end{figblock}



\subsection{وقتی $g$ یکنوا نیست: $Y=X^2$}
\label{sec:nonmono}

$x^2$ یکنوا نیست: $x$ و $-x$ هر دو به یک $y$ می‌روند، پس فرمول تغییر متغیر را نمی‌شود مستقیم به کار برد. اما روش \lr{CDF} اصلاً به یکنوایی نیاز نداشت. برای $y>0$:
\[ F_Y(y)=\P(X^2\le y)=\P(-\sqrt y\le X\le\sqrt y)=F_X(\sqrt y)-F_X(-\sqrt y), \]
و با مشتق‌گیری
\[ f_Y(y)=\frac{1}{2\sqrt y}\Bigl(f_X(\sqrt y)+f_X(-\sqrt y)\Bigr). \]
هر کدام از دو جمله سهمِ یکی از دو «شاخه» است: سهمِ $x>0$ و سهمِ $x<0$. هر شاخه به تنهایی یکنواست و فرمول تغییر متغیر را دارد؛ سهم‌ها جمع می‌شوند، درست مثل حالت گسسته که احتمالِ چند $x$ روی یک $y$ جمع می‌شد.

\begin{example}[label=ex:chi2]{مربعِ گاوسی استاندارد}
اگر $X\sim\mathcal N(0,1)$، چگالی $Y=X^2$ را بیابید.
\begin{solution}
با $f_X=\varphi$ و تقارن $\varphi(-z)=\varphi(z)$، برای $y>0$:
\[ f_Y(y)=\frac{1}{2\sqrt y}\cdot2\varphi(\sqrt y)=\frac{1}{\sqrt{2\pi y}}\,e^{-y/2}. \]
\end{solution}
\end{example}

%% ------------------------------------------------------------
\section{ساختن هر توزیعی از توزیع یکنواخت}
\label{sec:inverse}
%% ------------------------------------------------------------

تا اینجا $X$ را می‌دانستیم و دنبال توزیعِ $g(X)$ بودیم. حالا سؤال را برعکس می‌کنیم. فرض کنید فقط یک مولد عدد تصادفی داریم که $U\sim\mathrm{Unif}(0,1)$ می‌دهد؛ چیزی مثل \lr{\texttt{random()}} در پایتون یا \lr{\texttt{Math.random()}} در جاوااسکریپت. برای شبیه‌سازیِ یک صف، زمان‌های نمایی لازم داریم؛ برای شبیه‌سازی نویز، گاوسی؛ برای شبیه‌سازی تعداد تلاش تا موفقیت، هندسی. \key{چه تابعی $g$ را انتخاب کنیم که $g(U)$ دقیقاً \lr{CDF} دلخواهِ $F$ را داشته باشد؟}

\subsection{اول یک مشاهده: $F(X)$ همیشه یکنواخت است}

نمره‌های یک آزمون سراسری را در نظر بگیرید و فرض کنید نمره، $X$، توزیع پیوسته با \lr{CDF} برابر $F$ دارد. $F(x)$ یعنی چه کسری از شرکت‌کننده‌ها نمره‌شان حداکثر $x$ است؛ همان چیزی که به آن \emph{صدک} می‌گوییم. حالا یک شرکت‌کننده را تصادفی انتخاب کنید و صدکش، $F(X)$، را نگاه کنید. احتمال اینکه صدکش زیر $0.3$ باشد چقدر است؟ دقیقاً $0.3$، چون معنای صدکِ $0.3$ همین است که ۳۰ درصد افراد زیر آن‌اند. برای هر عدد دیگری هم همین استدلال کار می‌کند. پس صدکِ یک آدمِ تصادفی یکنواخت است، فرقی هم نمی‌کند نمره‌ها چه توزیعی داشته باشند.

\begin{proposition}[label=prop:pit]{یکنواخت بودنِ $F(X)$}
اگر $X$ پیوسته باشد و \lr{CDF} آن، $F$، روی برد $X$ اکیداً صعودی باشد، آنگاه $F(X)\sim\mathrm{Unif}(0,1)$.
\end{proposition}

\begin{proof}
$F(X)$ همیشه در $[0,1]$ است. برای $0<u<1$، $F$ در $u$ وارون دارد و چون $F$ صعودی است، $F(X)\le u$ همان $X\le F^{-1}(u)$ است. پس
\[ \P\bigl(F(X)\le u\bigr)=\P\bigl(X\le F^{-1}(u)\bigr)=F\bigl(F^{-1}(u)\bigr)=u, \]
که \lr{CDF} توزیع $\mathrm{Unif}(0,1)$ است.
\end{proof}

پیوسته بودن اینجا لازم است. اگر $X\sim\mathrm{Bern}(\tfrac12)$ باشد، $F(0)=\tfrac12$ و $F(1)=1$؛ پس $F(X)$ فقط دو مقدار $\tfrac12$ و $1$ را می‌گیرد و یکنواخت نیست.

\subsection{برعکس: $F^{-1}(U)$ توزیع $F$ را دارد}

$F$ متغیر $X$ را به یک عدد یکنواخت تبدیل کرد. پس طبیعی است که $F^{-1}$ یک عدد یکنواخت را به $X$ برگرداند.

\begin{theorem}[label=thm:inverse]{روش \lr{CDF} وارون}
فرض کنید $F$ یک \lr{CDF} پیوسته باشد که روی بازه‌ای که در آن $0<F<1$ است اکیداً صعودی است، و $U\sim\mathrm{Unif}(0,1)$. آنگاه \lr{CDF} متغیر $X=F^{-1}(U)$ دقیقاً $F$ است.
\end{theorem}

\begin{proof}
چون $F$ صعودی است، $F^{-1}(U)\le x$ همان $U\le F(x)$ است (کافی است $F$ را روی دو طرف اثر دهیم). پس
\[ \P(X\le x)=\P\bigl(F^{-1}(U)\le x\bigr)=\P\bigl(U\le F(x)\bigr)=F(x), \]
چون برای $0\le u\le1$، $\P(U\le u)=u$.
\end{proof}

اثبات این روش کوتاه است، اما ارزش دارد ببینیم که \emph{چرا} به درستی کار می‌کند. ابتدا نمودار تابع $F$ را رسم کنید. سپس یک عدد تصادفی یکنواخت روی محور عمودی انتخاب کنید، به صورت افقی جلو بروید تا به منحنی نمودار برخورد کنید، و از همان نقطه مستقیم به سمت پایین بیایید؛ عددی که روی محور افقی به آن می‌رسید، یک نمونه از متغیر $X$ است.

\begin{figblock}
\begin{tikzpicture}[x=2.2cm, y=4.2cm]
  \draw[-{Stealth}, graytext] (-0.05,0) -- (3.75,0) node[right] {$x$};
  \draw[-{Stealth}, graytext] (0,-0.03) -- (0,1.13) node[above] {$u=F(x)$};
  \draw[graytext, dashed] (0,1) -- (3.6,1) node[right, font=\scriptsize] {$1$};
  \draw[setA, very thick, domain=0:3.6, samples=80] plot (\x,{1-exp(-\x)});
  \foreach \uu/\xx/\lab/\dy in {0.1/0.105/0.11/0, 0.3/0.357/0.36/1, 0.5/0.693/0.69/0, 0.7/1.204/1.20/1, 0.9/2.303/2.30/0} {
    \draw[accent, dashed] (0,\uu) -- (\xx,\uu);
    \draw[accent, -{Stealth}] (\xx,\uu) -- (\xx,0.008);
    \fill[accent] (0,\uu) circle (1.6pt);
    \node[left, font=\scriptsize] at (0,\uu) {$\uu$};
    \node[below, font=\scriptsize] at (\xx,-0.01-0.06*\dy) {$\lab$};
  }
\end{tikzpicture}

\figcaption[fig:inverse]{روش تبدیل معکوس برای توزیع نمایی $\mathrm{Exp}(1)$ با فرمول $F(x)=1-e^{-x}$. پنج عدد که با فاصله‌های مساوی روی محور عمودی انتخاب شده‌اند، با فاصله‌های نامساوی روی محور افقی فرود می‌آیند: در نزدیکی صفر که شیب نمودار $F$ تند است، نمونه‌ها به هم نزدیک می‌افتند و در انتهای نمودار (دُم توزیع) که $F$ تقریباً صاف و افقی می‌شود، نمونه‌ها از هم دور می‌شوند.}
\end{figblock}

هر جا که شیب نمودار $F$ تند باشد (یعنی چگالی احتمال یا همان مشتق تابع، $f=F'$، بزرگ باشد)، یک بازهٔ پهن از مقادیر محور عمودی ($u$) روی یک بازهٔ باریک و کوچک از محور افقی ($x$) متمرکز می‌شود؛ در نتیجه، نمونه‌های زیادی در آن ناحیه جمع می‌شوند. برعکس، هر جا که نمودار $F$ صاف و کم‌شیب باشد، یک بازهٔ باریک از محور عمودی ($u$) روی بازهٔ پهنی از محور افقی ($x$) پخش می‌شود و تعداد نمونه‌ها در آنجا کم خواهد بود.

به بیان دقیق‌تر: نمونهٔ نهایی زمانی در بازهٔ کوچک $[x,\,x+dx]$ قرار می‌گیرد که عدد تصادفی $u$ در بازهٔ $[F(x),\,F(x+dx)]$ افتاده باشد. احتمال رخ دادن این اتفاق برابر است با طول این بازه، یعنی تقریباً $f(x)\,dx$؛ و این دقیقاً همان نتیجه‌ای است که به دنبالش بودیم.

\begin{keyformula}
\begin{enumerate}[leftmargin=1.4em, itemsep=4pt]
  \item \lr{CDF} توزیع موردنظر، $F$، را بنویسید.
  \item معادلهٔ $u=F(x)$ را برای $x$ حل کنید تا $x=F^{-1}(u)$ به دست آید.
  \item هر بار یک عدد یکنواخت $u$ بگیرید و $F^{-1}(u)$ را برگردانید.
\end{enumerate}
\end{keyformula}

\begin{example}[label=ex:inv-unif]{یکنواخت روی $[a,b]$}
از $U\sim\mathrm{Unif}(0,1)$ یک نمونه از $\mathrm{Unif}(a,b)$ بسازید.
\begin{solution}
$F(x)=\frac{x-a}{b-a}$ برای $a\le x\le b$. از $u=\frac{x-a}{b-a}$ نتیجه می‌شود $x=a+(b-a)u$. پس $X=a+(b-a)U$: بازهٔ $[0,1]$ را $b-a$ برابر می‌کشیم و $a$ واحد جابه‌جا می‌کنیم. همان چیزی است که بدون هیچ قضیه‌ای هم حدس می‌زدیم، و یک تبدیل خطی با ضریب مثبت است (بخش \ref{sec:linear}).
\end{solution}
\end{example}

\begin{example}[label=ex:inv-exp]{نمایی}
از $U\sim\mathrm{Unif}(0,1)$ یک نمونه از $\mathrm{Exp}(\lambda)$ بسازید. اگر $\lambda=2$ و مولد عدد $u=0.3$ را داده باشد، نمونه چیست؟
\begin{solution}
$u=1-e^{-\lambda x}$ را حل می‌کنیم: $e^{-\lambda x}=1-u$، پس
\[ x=F^{-1}(u)=-\frac{\ln(1-u)}{\lambda}. \]
برای $\lambda=2$ و $u=0.3$: \ $x=-\frac{\ln0.7}{2}\approx\frac{0.357}{2}\approx0.178$.

یک ساده‌سازی هم ممکن است. اگر $U$ یکنواخت باشد، $1-U$ هم یکنواخت است: ${\P(1-U\le u)=\P(U\ge1-u)=u}$. پس $X=-\frac{\ln U}{\lambda}$ هم $\mathrm{Exp}(\lambda)$ است. البته در عمل حواستان باشد: \lr{\texttt{random()}} در پایتون عددی در $[0,1)$ می‌دهد و ممکن است دقیقاً صفر برگرداند، و $\ln0$ تعریف نشده است. برای همین نسخهٔ ${-\ln(1-u)/\lambda}$ امن‌تر است:
\begin{center}
\lr{\texttt{x = -math.log(1 - random.random()) / lam}}
\end{center}
با این یک خط می‌شود صف بانک را شبیه‌سازی کرد: زمان بین ورود مشتری‌ها را این‌طور بسازید و پشت سر هم جمع کنید.
\end{solution}
\end{example}

\begin{example}[label=ex:inv-cauchy]{کوشی، و دُمِ سنگین در یک فرمول}
از $U\sim\mathrm{Unif}(0,1)$ یک نمونه از کوشی استاندارد بسازید و نمونه‌ای را که با $u=0.99$ ساخته می‌شود با نمونهٔ گاوسی استاندارد با همان $u$ مقایسه کنید.
\begin{solution}
در جلسهٔ ۸ دیدیم $F(x)=\frac12+\frac1\pi\arctan x$. از $u=\frac12+\frac1\pi\arctan x$:
\[ x=\tan\Bigl(\pi\bigl(u-\tfrac12\bigr)\Bigr). \]
برای $u=0.99$: \ $x=\tan(0.49\pi)\approx31.8$، در حالی که برای گاوسی استاندارد $\Phi^{-1}(0.99)\approx2.33$. وقتی $u$ به ۰ یا ۱ نزدیک می‌شود، زاویه به $\mp\frac\pi2$ نزدیک می‌شود و $\tan$ منفجر می‌شود. \lr{heavy tail} کوشی که در جلسهٔ ۸ با انتگرال دیدیم، اینجا مستقیم در فرمول پیداست.
\end{solution}
\end{example}



\subsection{حالت گسسته}

برای متغیر گسسته، $F$ پله‌ای است: جاهایی صاف است و جاهایی می‌پرد؛ پس وارونِ معمولی ندارد. ولی تصویرِ «از محور عمودی به محور افقی» هنوز کار می‌کند. بازهٔ $[0,1]$ را به تکه‌های پشت سر هم ببرید که طول هر تکه احتمالِ یکی از مقدارها باشد. اگر $U$ در تکهٔ $k$-ام افتاد، $x_k$ را برگردانید. احتمال افتادن $U$ در یک تکه برابر طول آن تکه است، یعنی همان $p_X(x_k)$.

\begin{figblock}
\begin{tikzpicture}[x=10cm, y=1cm]
  \fill[setA!35] (0,0) rectangle (0.2,0.7);
  \fill[accent!35] (0.2,0) rectangle (0.7,0.7);
  \fill[setC!35] (0.7,0) rectangle (1,0.7);
  \draw[graytext] (0,0) rectangle (1,0.7);
  \draw[graytext] (0.2,0) -- (0.2,0.7);
  \draw[graytext] (0.7,0) -- (0.7,0.7);
  \node[font=\small] at (0.1,0.35) {$X=1$};
  \node[font=\small] at (0.45,0.35) {$X=2$};
  \node[font=\small] at (0.85,0.35) {$X=3$};
  \foreach \tk in {0,0.2,0.7,1} \node[below, font=\scriptsize] at (\tk,0) {$\tk$};
  \draw[-{Stealth}, very thick, excolor] (0.65,1.45) -- (0.65,0.75);
  \node[above, font=\small, excolor] at (0.65,1.45) {$U=0.65$};
\end{tikzpicture}

\figcaption[fig:discrete-inverse]{ساختن متغیری با $p_X(1)=0.2$، $p_X(2)=0.5$ و $p_X(3)=0.3$. بازهٔ $[0,1]$ به سه تکه به همین طول‌ها بریده شده؛ $U=0.65$ در تکهٔ دوم افتاده، پس نمونه $X=2$ است.}
\end{figblock}

به زبان \lr{CDF}: \key{$X$ کوچک‌ترین $x$ای است که $F(x)\ge U$.} در شکل بالا، $F(1)=0.2$ و $F(2)=0.7$ و $F(3)=1$؛ برای $U=0.65$ کوچک‌ترین $x$ با $F(x)\ge0.65$ برابر ۲ است. چند نمونهٔ آشنا:
\begin{itemize}
  \item \textbf{برنولی:} $X=\ind\{U\le p\}$. یک تکه به طول $p$ برای ۱ و بقیه برای ۰.
  \item \textbf{تاس:} شش تکهٔ هم‌طول $\frac16$؛ پس $X=\lceil6U\rceil$. مثلاً $U=0.42$ می‌دهد $\lceil2.52\rceil=3$.
  \item \textbf{هندسی:} $F(k)=1-(1-p)^k$. کوچک‌ترین $k$ با $1-(1-p)^k\ge U$ را می‌خواهیم:
  \[ (1-p)^k\le1-U\iff k\ln(1-p)\le\ln(1-U)\iff k\ge\frac{\ln(1-U)}{\ln(1-p)}. \]
  در قدم آخر بر $\ln(1-p)$ تقسیم کردیم که منفی است، پس جهت نامساوی عوض شد (باز هم همان نکتهٔ علامت!). نتیجه: $X=\Bigl\lceil\frac{\ln(1-U)}{\ln(1-p)}\Bigr\rceil$. برای $p=\frac16$ (منتظر شش ماندن در تاس) و $U=0.5$: \ $\frac{\ln0.5}{\ln(5/6)}\approx\frac{-0.693}{-0.182}\approx3.80$، پس $X=4$.
\end{itemize}

فرمول هندسی را با نمایی مقایسه کنید: $-\frac{\ln(1-U)}{\lambda}$ با $\lambda=-\ln(1-p)$ دقیقاً همان کسرِ داخلِ $\lceil\ \rceil$ است. یعنی \emph{هندسی همان نمایی است که به بالا گرد شده}. در جلسهٔ ۸ نمایی را «هندسیِ پیوسته» نامیدیم؛ اینجا همان رابطه از طرف دیگر دیده می‌شود.

%% ------------------------------------------------------------
\section*{جمع‌بندی جلسهٔ دهم}
\addcontentsline{toc}{section}{خلاصه و تمرین‌ها}
%% ------------------------------------------------------------

\begin{summary}
\begin{itemize}
  \item \textbf{ایدهٔ اصلی:} $g(X)$ هم متغیر تصادفی است و $\P(g(X)\in B)=\P(X\in g^{-1}(B))$. در حالت گسسته $p_Y(y)=\sum_{x:\,g(x)=y}p_X(x)$.
  \item \textbf{روش \lr{CDF}:} برد $Y$ را پیدا کنید، $F_Y(y)=\P(g(X)\le y)$ را بر حسب $F_X$ بنویسید، و مشتق بگیرید. این روش برای هر $g$ای کار می‌کند.
  \item \textbf{تبدیل خطی:} $f_{aX+b}(y)=\frac1{|a|}f_X(\frac{y-b}{a})$. برای $a<0$ جهت نامساوی عوض می‌شود و $F_Y(y)=1-F_X(\frac{y-b}{a})$. اگر $X\sim\mathcal N(\mu,\sigma^2)$ آنگاه $aX+b\sim\mathcal N(a\mu+b,\ a^2\sigma^2)$.
  \item \textbf{تابع یکنوا:} $f_Y(y)=f_X(h(y))\,|h'(y)|$ با $h=g^{-1}$؛ یعنی $f_Y|dy|=f_X|dx|$. اگر $g$ یکنوا نیست، سهمِ شاخه‌ها جمع می‌شود.
  \item \textbf{رادیکال و نمایی:} $f_{\sqrt X}(y)=2y\,f_X(y^2)$ (یکنواخت $\to$ چگالی $2y$؛ نمایی $\to$ ریلی). $f_{e^X}(y)=\frac1y f_X(\ln y)$ (گاوسی $\to$ لگ‌نرمال).
  \item \textbf{\lr{CDF} وارون:} $F(X)\sim\mathrm{Unif}(0,1)$ و $F^{-1}(U)$ توزیع $F$ را دارد. نمایی: $-\frac{\ln(1-U)}{\lambda}$؛ کوشی: $\tan(\pi(U-\frac12))$؛ گسسته: $[0,1]$ را به تکه‌هایی به طول احتمال‌ها ببرید.
\end{itemize}
\end{summary}

\subsection*{تمرین‌ها و پاسخ تشریحی}

برای هر تمرین اول خودتان تلاش کنید و بعد پاسخ را بخوانید. در همهٔ پاسخ‌ها همان سه قدم روش \lr{CDF} را دنبال کرده‌ایم: برد، ترجمه، مشتق.

\begin{exercise}
اگر $X\sim\mathrm{Unif}(0,1)$، توزیع $Y=-\ln X$ را بیابید. این نتیجه را با مثال \ref{ex:inv-exp} مقایسه کنید.
\begin{solution}
\textbf{برد.} $X$ بین ۰ و ۱ است، پس $\ln X$ منفی است و $Y=-\ln X$ مثبت. وقتی $X$ به ۱ نزدیک می‌شود $Y$ به صفر می‌رسد و وقتی $X$ به صفر نزدیک می‌شود $Y$ به بی‌نهایت می‌رود. پس برد $Y$ بازهٔ $(0,\infty)$ است و برای $y\le0$، $F_Y(y)=0$.

\textbf{ترجمه.} برای $y>0$ نامساوی $-\ln X\le y$ را برای $X$ حل می‌کنیم. اول دو طرف را در $-1$ ضرب می‌کنیم (جهت عوض می‌شود)، بعد از دو طرف نمایی می‌گیریم ($e^x$ صعودی است و جهت را عوض نمی‌کند):
\[ -\ln X\le y\iff \ln X\ge-y\iff X\ge e^{-y}. \]
چون $0<e^{-y}<1$ و برای متغیر یکنواخت روی $(0,1)$ داریم $\P(X\ge t)=1-t$:
\[ F_Y(y)=\P(X\ge e^{-y})=1-e^{-y}. \]

\textbf{مشتق.} $f_Y(y)=e^{-y}$ برای $y>0$. این دقیقاً چگالی $\mathrm{Exp}(1)$ است.

\textbf{مقایسه.} در مثال \ref{ex:inv-exp} گفتیم $-\frac{\ln U}{\lambda}$ نمونه‌ای از $\mathrm{Exp}(\lambda)$ است. اینجا همان ادعا را برای $\lambda=1$ مستقیم و با روش \lr{CDF} ثابت کردیم.
\end{solution}
\end{exercise}

\begin{exercise}
اگر $X\sim\mathrm{Exp}(\lambda)$ و $c>0$، نشان دهید $cX\sim\mathrm{Exp}(\lambda/c)$. اگر زمان بر حسب دقیقه $\mathrm{Exp}(2)$ باشد، همان زمان بر حسب ثانیه چه توزیعی دارد؟
\begin{solution}
\textbf{برد.} $X\ge0$ و $c>0$، پس $Y=cX\ge0$. برای $y<0$، $F_Y(y)=0$.

\textbf{ترجمه.} برای $y\ge0$، چون $c$ مثبت است جهت نامساوی عوض نمی‌شود:
\[ F_Y(y)=\P(cX\le y)=\P\Bigl(X\le\frac yc\Bigr)=1-e^{-\lambda y/c}. \]
(از $F_X(x)=1-e^{-\lambda x}$ در جلسهٔ ۸ استفاده کردیم.)

\textbf{مشتق.} $f_Y(y)=\frac{\lambda}{c}\,e^{-(\lambda/c)y}$ برای $y\ge0$. این چگالی نمایی با پارامتر $\frac\lambda c$ است؛ پس $cX\sim\mathrm{Exp}(\lambda/c)$. همین را از گزارهٔ \ref{prop:linear} با $a=c$ و $b=0$ هم می‌شد گرفت.

\textbf{دقیقه به ثانیه.} اگر $X$ زمان بر حسب دقیقه باشد، همان زمان بر حسب ثانیه $Y=60X$ است. با $c=60$ و $\lambda=2$:
\[ Y\sim\mathrm{Exp}\Bigl(\frac{2}{60}\Bigr)=\mathrm{Exp}\Bigl(\frac1{30}\Bigr). \]
معنی‌اش هم روشن است: $\lambda$ «نرخ» است؛ ۲ پیشامد در هر دقیقه یعنی $\frac{2}{60}=\frac1{30}$ پیشامد در هر ثانیه. عوض کردن واحد، توزیع را از خانوادهٔ نمایی بیرون نمی‌برد و فقط نرخ را تغییر می‌دهد.
\end{solution}
\end{exercise}

\begin{exercise}
اگر $X\sim\mathcal N(0,1)$، چگالی $Y=|X|$ را بیابید.
\begin{solution}
\textbf{برد.} قدر مطلق منفی نمی‌شود، پس $Y\ge0$ و برای $y<0$، $F_Y(y)=0$.

\textbf{ترجمه.} $|x|$ یکنوا نیست ($x$ و $-x$ هر دو به یک عدد می‌روند)، پس فرمول تغییر متغیر را مستقیم به کار نمی‌بریم و از \lr{CDF} می‌رویم. برای $y\ge0$:
\[ F_Y(y)=\P(|X|\le y)=\P(-y\le X\le y)=\Phi(y)-\Phi(-y). \]
با تقارن $\Phi(-y)=1-\Phi(y)$ (جلسهٔ ۸):
\[ F_Y(y)=2\Phi(y)-1. \]

\textbf{مشتق.} چون $\Phi'=\varphi$:
\[ f_Y(y)=2\varphi(y)=\frac{2}{\sqrt{2\pi}}e^{-y^2/2}=\sqrt{\frac2\pi}\,e^{-y^2/2},\qquad y\ge0. \]

\textbf{شهود.} نیمهٔ منفیِ زنگوله روی نیمهٔ مثبت تا می‌خورد: هر $y>0$ هم از $x=y$ سهم می‌گیرد و هم از $x=-y$، و این دو سهم به خاطر تقارن برابرند؛ برای همین چگالی دو برابر شد. \textbf{چک:} $\int_0^\infty2\varphi(y)\,dy=2\times\frac12=1$.
\end{solution}
\end{exercise}

\begin{exercise}
اگر $X\sim\mathrm{Unif}(0,1)$، \lr{CDF} و چگالی $Y=e^X$ را بیابید.
\begin{solution}
\textbf{برد.} $X$ بین ۰ و ۱ است و $e^x$ صعودی است؛ پس $Y=e^X$ بین $e^0=1$ و $e^1=e$ است. در نتیجه برای $y<1$، $F_Y(y)=0$ و برای $y>e$، $F_Y(y)=1$.

\textbf{ترجمه.} برای $1\le y\le e$، از دو طرف $e^X\le y$ لگاریتم می‌گیریم ($\ln$ صعودی است و جهت را عوض نمی‌کند):
\[ F_Y(y)=\P(e^X\le y)=\P(X\le\ln y)=\ln y, \]
چون $0\le\ln y\le1$ و برای متغیر یکنواخت روی $(0,1)$، $\P(X\le t)=t$.

\textbf{مشتق.}
\[ F_Y(y)=\begin{cases}0 & y<1\\ \ln y & 1\le y\le e\\ 1 & y>e\end{cases}\qquad\qquad f_Y(y)=\begin{cases}\dfrac1y & 1\le y\le e\\[4pt] 0 & \text{\rl{در غیر این صورت}}\end{cases} \]

\textbf{چک.} $\int_1^e\frac{dy}{y}=\ln e-\ln1=1$. \lr{CDF} هم در $y=1$ از ۰ و در $y=e$ به ۱ پیوسته وصل می‌شود. چگالی نزدیک ۱ بزرگ‌تر است، چون $e^x$ نزدیک $x=0$ کُندتر از نزدیک $x=1$ رشد می‌کند و احتمال آنجا کمتر کش می‌آید.
\end{solution}
\end{exercise}

\begin{exercise}\hard
اگر $X\sim\mathrm{Unif}(-1,2)$، چگالی $Y=X^2$ را بیابید. (راهنمایی: بازه‌های $0<y\le1$ و $1<y\le4$ را جدا در نظر بگیرید.)
\begin{solution}
\textbf{برد.} $f_X(x)=\frac13$ روی $[-1,2]$. بزرگ‌ترین مقدار $X^2$ در $x=2$ است و برابر ۴ است و کوچک‌ترینش در $x=0$ و برابر ۰. پس برد $Y$ بازهٔ $[0,4]$ است.

\textbf{ترجمه.} برای $0<y\le4$:
\[ F_Y(y)=\P(X^2\le y)=\P(-\sqrt y\le X\le\sqrt y). \]
حالا باید دقت کنیم که بازهٔ $[-\sqrt y,\sqrt y]$ چقدرش داخل $[-1,2]$ است، یعنی جایی که $X$ واقعاً می‌تواند باشد. دو حالت پیش می‌آید:
\begin{itemize}
  \item \textbf{اگر $0<y\le1$:} آنگاه $\sqrt y\le1$؛ پس هر دو سرِ بازه، $-\sqrt y$ و $\sqrt y$، داخل $[-1,2]$ هستند. طول بازه $2\sqrt y$ است و
  \[ F_Y(y)=\frac{2\sqrt y}{3}. \]
  \item \textbf{اگر $1<y\le4$:} آنگاه $-\sqrt y<-1$، یعنی سرِ چپِ بازه بیرون از برد $X$ است. $X$ هیچ‌وقت کمتر از $-1$ نمی‌شود، پس عملاً بازهٔ $[-1,\sqrt y]$ را داریم با طول $\sqrt y+1$:
  \[ F_Y(y)=\frac{\sqrt y+1}{3}. \]
\end{itemize}

\textbf{مشتق.} با $(\sqrt y)'=\frac{1}{2\sqrt y}$:
\[ f_Y(y)=\begin{cases}\dfrac{2}{3}\cdot\dfrac{1}{2\sqrt y}=\dfrac{1}{3\sqrt y} & 0<y\le1\\[10pt] \dfrac13\cdot\dfrac{1}{2\sqrt y}=\dfrac{1}{6\sqrt y} & 1<y\le4\end{cases} \]

\textbf{چرا در $y=1$ چگالی نصف می‌شود؟} برای $y\le1$ هر $y$ از دو شاخه سهم می‌گیرد: $x=\sqrt y$ و $x=-\sqrt y$. برای $y>1$ شاخهٔ منفی دیگر وجود ندارد، چون $X$ به کمتر از $-1$ نمی‌رسد؛ فقط یک شاخه می‌ماند و چگالی نصف می‌شود.

\textbf{چک.} در $y=1$ هر دو فرمول \lr{CDF} عدد $\frac23$ را می‌دهند (پیوسته است) و در $y=4$ مقدار $\frac{2+1}{3}=1$ است. همچنین
\[ \int_0^1\frac{dy}{3\sqrt y}+\int_1^4\frac{dy}{6\sqrt y}=\frac23\bigl[\sqrt y\bigr]_0^1+\frac13\bigl[\sqrt y\bigr]_1^4=\frac23+\frac13=1. \]
\end{solution}
\end{exercise}

\begin{exercise}
با روش \lr{CDF} وارون، از $U\sim\mathrm{Unif}(0,1)$ یک نمونه از لاپلاس استاندارد، با چگالی $\frac12e^{-|x|}$، بسازید.
\begin{solution}
\textbf{قدم ۱: \lr{CDF}.} چگالی دو تکه دارد، پس \lr{CDF} را هم دو تکه حساب می‌کنیم.\\
برای $x<0$ فقط تکهٔ چپ، $\frac12e^{t}$، را جمع می‌زنیم:
\[ F(x)=\int_{-\infty}^x\tfrac12e^{t}\,dt=\tfrac12e^{x}. \]
برای $x\ge0$ کل نیمهٔ چپ را داریم (که مساحتش $\frac12$ است) و از ۰ تا $x$ هم تکهٔ راست، $\frac12e^{-t}$، اضافه می‌شود:
\[ F(x)=\tfrac12+\int_0^x\tfrac12e^{-t}\,dt=\tfrac12+\tfrac12\bigl(1-e^{-x}\bigr)=1-\tfrac12e^{-x}. \]
دقت کنید $F(0)=\frac12$؛ یعنی $u<\frac12$ متناظر با $x$ منفی و $u\ge\frac12$ متناظر با $x$ نامنفی است.

\textbf{قدم ۲: حل $u=F(x)$.}
\begin{itemize}
  \item اگر $u<\frac12$: \ $u=\frac12e^{x}\ \Rightarrow\ e^x=2u\ \Rightarrow\ x=\ln(2u)$. (چون $2u<1$، این $x$ منفی است، همان‌طور که باید.)
  \item اگر $u\ge\frac12$: \ $u=1-\frac12e^{-x}\ \Rightarrow\ e^{-x}=2(1-u)\ \Rightarrow\ x=-\ln\bigl(2(1-u)\bigr)$. (چون $2(1-u)\le1$، این $x$ نامنفی است.)
\end{itemize}

\textbf{قدم ۳: نمونه‌گیری.}
\[ X=\begin{cases}\ln(2U) & U<\frac12\\[3pt] -\ln\bigl(2(1-U)\bigr) & U\ge\frac12\end{cases} \]
\textbf{چک با عدد.} $U=0.25$ می‌دهد $\ln0.5\approx-0.69$ و $U=0.75$ می‌دهد $-\ln0.5\approx+0.69$. دو $u$ که نسبت به $\frac12$ قرینه‌اند، دو $x$ قرینه می‌دهند؛ این با تقارن چگالی لاپلاس حول صفر جور است.
\end{solution}
\end{exercise}

\begin{exercise}
$X$ مقادیر $0$، $1$، $2$ و $3$ را با احتمال‌های $0.1$، $0.4$، $0.3$ و $0.2$ می‌گیرد. مولد عدد یکنواخت سه عدد $0.05$، $0.47$ و $0.93$ را داده است. با روش این جلسه سه نمونه از $X$ بسازید.
\begin{solution}
\textbf{قدم ۱: \lr{CDF}.} احتمال‌ها را پشت سر هم جمع می‌کنیم:
\[ F(0)=0.1,\qquad F(1)=0.1+0.4=0.5,\qquad F(2)=0.5+0.3=0.8,\qquad F(3)=0.8+0.2=1. \]

\textbf{قدم ۲: تکه‌ها.} بازهٔ $[0,1]$ با همین نقطه‌ها به چهار تکه بریده می‌شود:
\begin{center}
\renewcommand{\arraystretch}{1.25}
\begin{tabular}{c|cccc}
\rl{تکه} & $[0,\,0.1]$ & $(0.1,\,0.5]$ & $(0.5,\,0.8]$ & $(0.8,\,1]$\\\hline
\rl{طول} & $0.1$ & $0.4$ & $0.3$ & $0.2$\\
\rl{نمونه} & $X=0$ & $X=1$ & $X=2$ & $X=3$
\end{tabular}
\end{center}
طول هر تکه دقیقاً احتمالِ مقدارِ متناظرش است؛ برای همین روش درست کار می‌کند.

\textbf{قدم ۳: نمونه‌ها.} برای هر $u$ کوچک‌ترین $x$ با $F(x)\ge u$ را پیدا می‌کنیم:
\begin{itemize}
  \item $u=0.05$: \ $F(0)=0.1\ge0.05$، پس $X=0$.
  \item $u=0.47$: \ $F(0)=0.1<0.47$ ولی $F(1)=0.5\ge0.47$، پس $X=1$.
  \item $u=0.93$: \ $F(2)=0.8<0.93$ ولی $F(3)=1\ge0.93$، پس $X=3$.
\end{itemize}
پس سه نمونه $0$، $1$ و $3$ هستند.
\end{solution}
\end{exercise}
